% Pursuing Stacks (cotes 134-2 to 134-8): archive cote 134-5, site page 31 (PDF page 32), typescript page 218.
% typescript. Draft compared: ps4.tex:1480-1551.
% AI-transcribed by Claude Opus 5.5 (claude-opus-5-5, adaptive thinking, effort high), Message Batches API, batch ps2 (msgbatch_01NVwr1Yuo8qJf5rykeGBXPK), 2026-11-12; reviewed against both supplied scan halves; see CHECK_DECISIONS.tsv.
\scanpage{134-5/31}
\mspage{218}

tion on B,  there is \xout{such an interval} \inserted{a} separating $\underline{W}_B$-homotopy interval
$\underline{I} = (I,\hand{$\delta$}_0,\hand{$\delta$}_1)$ in $B^\wedge$.  The exactness properties of $i^*$  imply that
$i^*(\underline{I})$  is a separating interval in $A^\wedge$,  \xout{and} the asphericity condition on
i implies that moreover $i^*(I)$ is $\underline{W}$-aspheric,  hence $\underline{W}$-aspheric over $e_{A^\wedge}$
% A pen spacing slash separates the following A-hat from is; it is not a mathematical slash.
as $A^\wedge$ is  totally $\underline{W}$-aspheric,  qed.

To finish \xout{these} with the more formal properties of the notion of
$\underline{W}$-aspheric maps in (Cat),  let's give a list of the standard stability
conditions for this notion,  with respect notably to composition,  base
change,  and cartesian products:

\hand{\uline{\textcolor{black}{Proposition 3}}}.   a)  Consider two maps  \xout{AxxxxxB}
\[ A \overset{i}{\hand{\longrightarrow}} B \overset{j}{\hand{\longrightarrow}} C \]
in (Cat).   Then if  i and j are $\underline{W}$-aspheric,  so is ji.  If ji and i are
$\underline{W}$-aspheric,  so is j .  Any isomorphism in (Cat) is $\underline{W}$-aspheric.

b)   Let
% Hand-drawn arrows are marked blue; typed objects remain plain.
% A long pen curve runs beneath this square and partly over the next line; its function is unresolved.
\diagram{\begin{tikzcd}
A & A' \arrow[color=blue, l, "f"'] \arrow[color=blue, d, "i'"] \\
B \arrow[color=blue, from=u, "i"'] & B' \arrow[color=blue, l, "g"]
\end{tikzcd}}
be a cartesian square in (Cat),  assume i is $\underline{W}$-aspheric and g is \xout{xx}fibe-%
ring  (for instance an induction functor  $B_{/F} \hand{\to} B$ ,  with  F in $B^\wedge$).
Then i' is $\underline{W}$-aspheric  (and, of course,  f is equally fibering).

c)   Let
\[ i: A \hand{\to} B , \quad i' : A' \hand{\to} B' \]
be two W-aspheric maps in (Cat),  then
\[ i\mathrm{x}i' : A\mathrm{x}A' \hand{\to} B\mathrm{x}B' \]
is $\underline{W}$-aspheric.

\hand{\uline{\textcolor{black}{Proof}}}.  Property a) is formal,  in terms of the criterion (ii') of
prop. 1 .  Property c) follows formally from the criterion (iii),  and
the canonical isomorphism
\[ (A\mathrm{x}A')_{/b\mathrm{x}b'} \hand{\simeq} (A_{/b})\mathrm{x}(A'_{/b'}) \quad , \]
and the fact that a product of two $\underline{W}$-aspheric objects of (Cat) is again
$\underline{W}$-aspheric  (cf.  prop. of page 167,  making use of  \xout{Loc 3)} the localiza-%
tion condition  Loc 3)  on $\underline{W}$  in it's full generality).   Another proof
goes via  a) and b),  by reducing first (using a)) to the case when
either i or i' are identities,  and using the fact that a projection map
$C\mathrm{x}B \hand{\to} B$ in (Cat) is a fibration.

We are left with proving property b).  For this,  we note that the
asphericity condition (iii) on a map $i: A \hand{\to} B$  just means that \xout{the} for
\margin{31}
